
% =============================================================
%  Chapter 1 — Introduction
% =============================================================
\chapter{Introduction}
\label{ch:introduction}

% ----- §1.1  Broad context: LLMs generating code -----
\section{LLM-Based Code Generation}
\label{sec:intro-context}

Large language models have advanced rapidly from completing individual
functions~\cite{chen2021humaneval} to resolving real-world software engineering
issues across entire repositories~\cite{jimenez2024swebench}. Today, models
routinely generate multi-file backend applications that span frameworks,
databases, and authentication layers. \textsc{BaxBench}~\cite{vero2025baxbench},
the benchmark this thesis extends, evaluates this capability directly: given a
fixed backend task, it checks whether a model's generated application is both
functionally correct and free of exploitable vulnerabilities
(Section~\ref{sec:background-codegen} details its structure). Even the
best-performing models reach only around 62\% functional correctness, and
roughly half of the implementations that do pass remain exploitable, so code
generation is far from a solved problem at the level of a complete
application.

\textsc{BaxBench}, like the function- and repository-level benchmarks that
precede it, shares a fundamental limitation in evaluation methodology: it
executes each generated application in an isolated, single-machine container
and measures only whether the code passes functional and security tests. No
assessment is made of how the application would \emph{perform} under
realistic deployment conditions.


% ----- §1.2  Narrow the gap: performance under deployment -----
\section{Beyond Correctness: Performance Under Real Deployment}
\label{sec:intro-gap}

In production, backend applications do not run in isolation. They are
deployed on container orchestration platforms such as
Kubernetes~\cite{burns2016borg,kubernetes} (Chapter~\ref{ch:background}
introduces the relevant concepts), where multiple replicas share finite CPU,
memory, and network resources across a cluster of worker nodes. Achieving
high throughput under these conditions requires careful \emph{deployment
configuration}: decisions about replica counts, resource allocation, pod
placement, and database topology
(Section~\ref{sec:background-deployment-levers}).

A functionally correct application can still fail to serve production traffic
if these parameters are poorly chosen. Today, they are tuned manually by
operations engineers or handled by reactive autoscaling mechanisms such as the
Kubernetes Horizontal and Vertical Pod Autoscalers~\cite{kubernetes-hpa},
neither of which reasons proactively about the interaction between
application behaviour and infrastructure layout, and neither of which can fix
a bug in the application code itself.

Generating correct code is therefore a necessary but insufficient condition
for production readiness. \emph{Deployment configuration is a critical and
largely unexamined dimension of LLM-based software generation, and it does
not act in isolation from code correctness}: as this thesis's own results
later show (Chapter~\ref{ch:results}), some of the most consequential
deployment failures are in fact latent code bugs that only surface once an
application is deployed behind real cluster infrastructure. No existing
benchmark evaluates whether LLMs can reason about, and iteratively improve,
the deployment of the applications they generate, jointly with the code
those deployments run.

Evaluating this problem therefore requires more than executing an application
in an isolated container. Candidates must be observed on a multi-node
Kubernetes cluster, where replica placement, shared resource capacity, and
database topology can affect the performance that the benchmark measures.


% ----- §1.3  The gap stated explicitly -----
\section{The Research Gap: Deployment Optimisation as an Unexamined Dimension}
\label{sec:intro-missing}

Two lines of prior work come closest to this problem, and both stop short of
it. One optimises performance at the level of source code or system
internals: synthesising bespoke LLM serving runtimes~\cite{kamahori2025vibeserve},
entire key-value stores~\cite{liu2025jitsystems}, or repository-scale
performance patches~\cite{he2025sweperf}. All three iterate, correctness-gate,
and measure, but none touches how a fixed application is configured,
replicated, and placed on a cluster. The other targets deployment directly,
but not performance: Glia~\cite{hamadanian2025glia} combines LLM reasoning
with measured feedback to design GPU-cluster scheduling and routing
mechanisms, the closest prior system to this thesis's setting, but for
inference-serving mechanism design rather than arbitrary backend deployment;
\textsc{IaCGen}~\cite{zhang2026iacgen} and TerraFormer~\cite{jana2026terraformer}
iteratively regenerate infrastructure-as-code from deployment feedback, but
stop at \emph{deployability}, a binary signal, rather than the throughput or
latency the provisioned system can sustain under load; and learned
cluster-management systems such as Autopilot~\cite{rzadca2020autopilot},
FIRM~\cite{qiu2020firm}, and Decima~\cite{mao2019decima} optimise
resource-management policy from data, but hold the application and its
deployment manifest fixed and are not driven by an LLM.
Chapter~\ref{ch:related-work} surveys this space in full; the point for now
is that no existing system combines an LLM-controlled deployment
specification, an LLM-controlled application, and a runtime performance
objective in one iterative loop.

This gap leaves a fundamental question unanswered:
%
\begin{quote}
  \emph{Can an LLM agent iteratively optimise the deployment configuration of
  a backend application on a Kubernetes cluster, guided by runtime performance
  feedback, to maximise sustainable throughput under fixed infrastructure
  constraints?}
\end{quote}


% ----- §1.4  This work -----
\section{This Work}
\label{sec:intro-thiswork}

In this thesis, we extend \textsc{BaxBench} with a Kubernetes-based iterative
benchmarking framework that addresses this question.  Given a
\textsc{BaxBench} scenario and a target Kubernetes cluster, the framework
orchestrates a multi-phase refinement loop:
%
\begin{enumerate}
  \item A \textbf{baseline phase} in which the LLM agent generates both
        application code and an initial deployment specification, which are
        validated, deployed, and benchmarked under adaptive load.
  \item An \textbf{iterative refinement phase} in which the agent receives
        structured performance feedback, load-test statistics, per-endpoint
        latency distributions, resource utilisation diagnostics, and the
        previous deployment's goodput, and decides whether to improve the
        application code \emph{or} the deployment configuration, exclusively,
        before repeating the deploy--benchmark cycle.
\end{enumerate}
%
The deployment specification is expressed as a constrained YAML schema
(\texttt{spec.yaml}) that gives the agent control over the deployment levers
introduced in Section~\ref{sec:background-deployment-levers}, replica count,
resource allocation, placement, and database topology, while the framework
handles manifest rendering, Kubernetes resource management, and
orchestration.

Performance is measured using an adaptive load profile
(\texttt{k8s-explore-refine}) that automatically discovers the sustainable
throughput of each deployment.  Rather than fixing the request rate a priori,
the load controller ramps concurrent virtual users in steps until
\emph{goodput} (Section~\ref{sec:background-goodput}) stops growing
meaningfully.  This makes iteration scores directly comparable across
applications, deployment sizes, and refinement rounds without per-scenario
tuning.  Every LLM call is logged against its measured cost, so that model
comparisons in Chapter~\ref{ch:results} can be read alongside capability
rather than in place of it.


% ----- §1.5  Research questions -----
\section{Research Questions}
\label{sec:intro-rqs}

The central question above is broad enough that answering it in aggregate
would hide more than it reveals: whether refinement helps, which lever drives
that help, whether it depends on the model, and where it breaks down are
separate empirical claims. Chapter~\ref{ch:results} therefore decomposes it
into six questions, asked of the same dataset (three models, eight
\textsc{BaxBench} scenarios, up to three frameworks; Chapter~\ref{ch:setup}):
%
\begin{description}
  \item[RQ1] Does iterative refinement improve sustained goodput
    (Section~\ref{sec:results-rq1})?
  \item[RQ2] Which refinement lever, code or deployment, contributes more,
    and do the two specialise on different failure classes
    (Section~\ref{sec:results-rq2})?
  \item[RQ3] How do the evaluated models compare, in refinement outcome and
    in cost (Section~\ref{sec:results-rq3})?
  \item[RQ4] Does application-framework choice affect achievable goodput
    independently of refinement (Section~\ref{sec:results-rq4})?
  \item[RQ5] What deployment parameters does the agent converge to, and do
    they track the workload's actual characteristics
    (Section~\ref{sec:results-rq5})?
  \item[RQ6] What are the dominant failure modes, and what do individual
    trajectories reveal that aggregate statistics do not
    (Section~\ref{sec:results-rq6})?
\end{description}
%
Chapter~\ref{ch:discussion} interprets the six answers together, and
Chapter~\ref{ch:conclusion} returns to the overarching question directly.


% ----- §1.6  Contributions -----
\section{Contributions}
\label{sec:intro-contributions}

This thesis makes the following contributions:
%
\begin{enumerate}
  \item \textbf{An iterative benchmark framework} that extends
        \textsc{BaxBench} with Kubernetes deployment, adaptive load testing,
        and a multi-phase refinement loop, decision, then \emph{either} code
        \emph{or} spec regeneration, then deploy, then bench, then feedback,
        enabling the evaluation of LLM-driven deployment optimisation
        alongside code generation. A failure taxonomy and deterministic
        routing scheme (Section~\ref{sec:methods-failure}) keeps a single
        failed iteration from silently corrupting the next one, which is what
        makes the empirical results below attributable to a specific cause
        rather than an unexplained end-to-end number.

  \item \textbf{A constrained deployment specification} (\texttt{spec.yaml})
        that exposes a structured search space for the LLM agent, covering
        replica count, resource requests and limits, database topology, and pod
        placement, while the framework enforces cluster-capacity validation
        and deterministic manifest rendering.

  \item \textbf{An adaptive load profile} (\texttt{k8s-explore-refine}) that
        automatically finds the sustainable throughput of any deployment
        without per-application rate tuning, ramping load through an
        explore phase to locate an approximate overload boundary, then a
        refine phase that searches near it in efficiency-scaled steps,
        with overload and stall detection based on failure rate, tail
        latency, and goodput collapse thresholds.

  \item \textbf{An experimental evaluation} across three frontier-tier
        LLMs, eight \textsc{BaxBench} scenarios, and up to three
        application frameworks (72 scenario$\times$framework$\times$model
        cells in the full matrix, 63 with experiment data), tracking LLM
        cost as a first-class quantity alongside goodput. The evaluation
        shows that iterative refinement improves sustained goodput in
        essentially every cell with data (gains from just over $1\times$ to
        more than $75\times$ depending on the baseline), in several cases
        converting a non-functional deployment into a working one; that code
        and deployment refinement specialise in practice on different
        failure classes, with code refinement carrying the largest
        single-step recoveries and deployment refinement acting as a
        smaller, more model-dependent capacity knob; and that all three
        models reach a working baseline and complete their full refinement
        budget on every cell with data, so model differences show up in
        refinement behaviour, such as one model's deployment-refinement
        steps turning net negative on average, rather than in whether a
        baseline is reached at all (Chapter~\ref{ch:results}). Taken
        together, these findings identify entanglement with code
        correctness, not the deployment-reasoning task itself, as the
        largest current limitation of LLM-guided deployment optimisation.
\end{enumerate}


% ----- §1.7  Thesis outline -----
\section{Thesis Outline}
\label{sec:intro-outline}

The remainder of this thesis is organised as follows.
%
\begin{description}
  \item[Chapter~\ref{ch:background}] introduces the technical background this
    thesis assumes: \textsc{BaxBench} and the limits of single-container
    evaluation, Kubernetes and the deployment levers it exposes, and the
    vocabulary, correctness versus performance, throughput, latency, and
    goodput, used to evaluate deployments in later chapters.

  \item[Chapter~\ref{ch:related-work}] surveys related work on LLM code
    generation and evaluation, iterative self-refinement and agentic coding,
    verification-gated generation, deployment configuration and learned
    cluster management, and LLM-driven performance optimisation, and
    positions this thesis against it.

  \item[Chapter~\ref{ch:methods}] describes the system architecture, the
    iterative experiment loop and its individual phases, the adaptive load
    profile and goodput metric, failure handling, and the constrained
    deployment specification.

  \item[Chapter~\ref{ch:setup}] details the experimental setup, including
    cluster configuration, selected models, scenarios, and hyperparameters.

  \item[Chapter~\ref{ch:results}] presents the experimental results, organised
    around the six research questions of Section~\ref{sec:intro-rqs}, closing
    with a synthesis that connects them back to the overarching question of
    Section~\ref{sec:intro-missing}.

  \item[Chapter~\ref{ch:discussion}] interprets those results, discusses their
    limitations and threats to validity, and draws out practical implications
    for agentic software optimisation and Kubernetes deployment tuning.

  \item[Chapter~\ref{ch:conclusion}] answers the overarching research question
    directly, summarises what worked and what did not, and outlines
    directions for future work.
\end{description}
